% Generated draft: requires mathematical and source review.
\documentclass[11pt]{article}
\usepackage[margin=25mm]{geometry}
\usepackage{fontspec}
\setmainfont{Latin Modern Roman}
\usepackage{amsmath,amssymb,mathtools,unicode-math}
\setmathfont{Latin Modern Math}
\usepackage{enumitem,needspace,xurl,hyperref}
\setlength{\parindent}{0pt}
\setlength{\parskip}{5pt}
\setlength{\emergencystretch}{3em}
\newcommand{\N}{\mathbb{N}}
\newcommand{\Z}{\mathbb{Z}}
\newcommand{\Q}{\mathbb{Q}}
\newcommand{\R}{\mathbb{R}}
\newcommand{\C}{\mathbb{C}}
\newcommand{\F}{\mathbb{F}}
\newcommand{\A}{\mathbb{A}}
\newcommand{\PP}{\mathbb{P}}
\newcommand{\E}{\mathbb{E}}
\newcommand{\eps}{\varepsilon}
\newcommand{\dd}{\,\mathrm{d}}
\begin{document}
\section*{UnsolvedMath ID: 30001064}
\textbf{Source problem code:} catalogue entry
\subsection{Definitions and mathematical statement}
For a planar set \(P\subseteq\R^2\), let \(\mathcal F=(P+t_i)_{i\in I}\) be an indexed family of translates, with each copy assigned its own index. It is a \(k\)-fold covering of the plane when
\[
 \forall x\in\R^2,\qquad
 |\{i\in I:x\in P+t_i\}|\geq k.
\]
Thus \(k\)-fold means at least \(k\), not exactly \(k\). A decomposition into two coverings means a partition \(I=I_0\mathbin{\dot\cup}I_1\) such that
\[
 \forall a\in\{0,1\}\ \forall x\in\R^2\ \exists i\in I_a:
 x\in P+t_i.
\]
Equivalently, color each translate red or blue so every point lies in a translate of each color. The copies are partitioned; the sets themselves are not cut into pieces.

The cover-decomposability assertion for \(P\) is
\[
 \exists k=k(P)\geq2\ \forall\mathcal F\text{ a }k\text{-fold translate covering of }\R^2,
 \quad\exists\text{ such a two-color decomposition}.
\]
The question attributed to the Pach conjecture asks whether this holds for every planar convex set \(P\). The threshold may depend on \(P\); no universal threshold for all convex sets is requested. Only translations of the fixed set are allowed, without rotations or rescalings.

The record does not specify whether covering families must be locally finite, countable, or satisfy another multiplicity condition, or whether repeated translates are permitted; the indexed-family definition makes repetitions meaningful but needs agreement with the report's convention. Infinite families are essential for covering the whole plane by translates of a bounded set. It also does not restrict convex sets to bounded bodies with nonempty interior, or select open versus closed sets. Boundary membership follows the actual chosen set and can affect the covering property.

The supplied assessment dated 2026-08-21 reports cover-decomposability for open convex polygons and calls the unrestricted planar-convex-set conjecture unresolved. This partial result is not extended to closed polygons or arbitrary curved convex sets without the missing theorem hypotheses.
\subsection{Short English statement}
Does every fixed convex planar set have a multiplicity threshold above which any translate covering can be colored into two full coverings? The threshold depends on the set, and the family and boundary conventions need to match the source.
\subsection{Sources}
\begin{itemize}
\item Local UnsolvedMath registry, permanent ID \(30001064\), exact source code \texttt{catalogue entry}, proposed year 2008; no individual proposer is supplied. The assessment identifies the general question as the Pach conjecture.
\item \emph{Discrete Geometry (2009)}, Oberwolfach Reports, DOI 10.4171/owr/2008/44, pp. 2477--2556. \url{https://doi.org/10.4171/owr/2008/44}.
\item Supplied reference labelled ``Survey on multiple coverings'': \url{https://math.bme.hu/~zv/kivonatok/2015/kovacs_istvan_diplkiv.pdf}.
\item The supplied \texttt{corrected\_verified} note reports meaning-preserving semantic and notation cleanup without unsupported added content. The open label and assessment checked 2026-08-21 are supplied dated curation, not independently verified; the precise open-polygon theorem and current general status require review.
\end{itemize}
\end{document}
